% =====================================================================
% CAPITULO 9 --- ANALISE DE MALHAS / METODO DE MAXWELL --- V2
% Estrutura limpa: teoria curta + banco visual diversificado por nivel.
% As fontes dependentes sao inseridas nos niveis correspondentes, sem
% transformar o capitulo inteiro em variacoes do mesmo circuito.
% =====================================================================

\section{Método de Maxwell (Análise de Malhas)}

\begin{conceito}[Procedimento de análise]
Para circuitos \textbf{planares}, associe uma corrente a cada malha interna.
Neste material, a convenção gráfica é sempre
\textbf{corrente de malha no sentido horário}.
\begin{enumerate}[leftmargin=1.7em,itemsep=2pt,topsep=3pt]
\item Identifique as malhas independentes e adote $I_1,I_2,\ldots$ no sentido horário.
\item Em resistor exclusivo da malha $k$, use $RI_k$.
\item Em resistor compartilhado por $j$ e $k$, use
      $R(I_j-I_k)$ quando o percurso é feito na malha $j$.
\item Fontes de tensão entram diretamente na LKT, com sinal definido pela
      polaridade encontrada durante o percurso horário.
\item Fonte de corrente na periferia fixa diretamente uma corrente de malha.
\item Fonte de corrente em ramo compartilhado exige \textbf{supermalha} mais
      uma equação de restrição.
\item Resolva o sistema e converta correntes de malha em correntes reais de ramo.
\end{enumerate}
\end{conceito}

\begin{atencao}[Sentido das setas]
As setas desenhadas neste capítulo são \textbf{horárias}. Em TikZ/CircuitikZ,
um arco com ângulo final positivo é anti-horário; por isso os desenhos usam
ângulos finais negativos. Se uma corrente calculada resultar negativa, o
\emph{resultado} circula no sentido contrário à seta, mas a seta de referência
continua horária.
\end{atencao}

\legendaniveis

% Injecao dos bancos de fontes dependentes exatamente entre os niveis do banco
% principal. Isso preserva a classificacao N2/N3/N4.
\begingroup
\newcount\fvMaxAprof
\newcount\fvMaxDesafio
\fvMaxAprof=0
\fvMaxDesafio=0

\let\fvMaxAprofOrig\blocoaprofundamento
\renewcommand{\blocoaprofundamento}{%
  \advance\fvMaxAprof by 1\relax
  \ifnum\fvMaxAprof=1\relax
    % =====================================================================
% FONTES DEPENDENTES EM MALHAS --- NIVEL 2
% Questoes integradas ao bloco de Aplicacao do Capitulo 9.
% Correntes de malha desenhadas no sentido HORARIO.
% =====================================================================

\begin{questao}[2]{}
No circuito, a fonte de tensão dependente fornece $2{,}5I_1$ volts e atua na
malha 2. Determine $I_1$, $I_2$ e a corrente real no resistor central.

\circuitoq{%
\begin{circuitikz}[american,scale=.92,transform shape,line width=.8pt]
\draw (0,0) to[\fonteV,l^={$\SI{20}{\volt}$}] (0,3.2)
      to[R,l={$\SI{5}{\ohm}$}] (3.2,3.2) coordinate(t);
\draw (t) to[R,l={$\SI{5}{\ohm}$}] (3.2,0) coordinate(b) -- (0,0);
\draw (t) to[R,l={$\SI{5}{\ohm}$}] (6.4,3.2)
      to[cvsource,l={$2{,}5I_1$}] (6.4,0) -- (b);
\draw[->,loopcol,line width=1.2pt] (1.75,1.00) arc[start angle=0,end angle=-310,radius=.46cm];
\node[loopcol,fill=white,inner sep=1pt] at (1.75,1.00){\footnotesize$I_1$};
\draw[->,loopcol,line width=1.2pt] (4.75,1.00) arc[start angle=0,end angle=-310,radius=.46cm];
\node[loopcol,fill=white,inner sep=1pt] at (4.75,1.00){\footnotesize$I_2$};
\end{circuitikz}}

\resposta{$I_1=\SI{3.2}{\ampere}$; $I_2=\SI{2.4}{\ampere}$; $I_c=\SI{0.8}{\ampere}$}
\resolucao{Com correntes de malha horárias,
\[
10I_1-5I_2=20,
\qquad
-5I_1+10I_2=2{,}5I_1.
\]
A segunda equação dá $I_2=0{,}75I_1$. Substituindo na primeira,
$I_1=\SI{3.2}{\ampere}$ e $I_2=\SI{2.4}{\ampere}$. No ramo comum,
$I_c=I_1-I_2=\SI{0.8}{\ampere}$.}
\end{questao}

\begin{questao}[2]{}
A fonte de corrente dependente está na periferia da malha 2 e impõe
$I_2=0{,}5I_1$. Determine $I_1$, $I_2$ e o módulo da tensão sobre essa fonte.

\circuitoq{%
\begin{circuitikz}[american,scale=.92,transform shape,line width=.8pt]
\draw (0,0) to[\fonteV,l^={$\SI{30}{\volt}$}] (0,3.2)
      to[R,l={$\SI{5}{\ohm}$}] (3.2,3.2) coordinate(t);
\draw (t) to[R,l={$\SI{5}{\ohm}$}] (3.2,0) coordinate(b) -- (0,0);
\draw (t) to[R,l={$\SI{10}{\ohm}$}] (6.4,3.2)
      to[cisource,l={$0{,}5I_1$},invert] (6.4,0) -- (b);
\draw[->,loopcol,line width=1.2pt] (1.75,1.00) arc[start angle=0,end angle=-310,radius=.46cm];
\node[loopcol,fill=white,inner sep=1pt] at (1.75,1.00){\footnotesize$I_1$};
\draw[->,loopcol,line width=1.2pt] (4.75,1.00) arc[start angle=0,end angle=-310,radius=.46cm];
\node[loopcol,fill=white,inner sep=1pt] at (4.75,1.00){\footnotesize$I_2$};
\end{circuitikz}}

\resposta{$I_1=\SI{4}{\ampere}$; $I_2=\SI{2}{\ampere}$; $|v_s|=\SI{10}{\volt}$}
\resolucao{Por estar na periferia, a fonte fixa diretamente
$I_2=0{,}5I_1$. Na malha 1,
\[
10I_1-5I_2=30.
\]
Logo $10I_1-2{,}5I_1=30$, de onde $I_1=\SI{4}{\ampere}$ e
$I_2=\SI{2}{\ampere}$. Na malha 2, as quedas resistivas totalizam
$10(2)+5(2-4)=\SI{10}{\volt}$ em módulo; a fonte de corrente sustenta essa tensão.}
\end{questao}
%
  \fi
  \fvMaxAprofOrig
}

\let\fvMaxDesafioOrig\blocodesafio
\renewcommand{\blocodesafio}{%
  \advance\fvMaxDesafio by 1\relax
  \ifnum\fvMaxDesafio=1\relax
    % =====================================================================
% FONTES DEPENDENTES EM MALHAS --- NIVEL 3
% Questoes integradas ao bloco de Aprofundamento do Capitulo 9.
% Correntes de malha desenhadas no sentido HORARIO.
% =====================================================================

\begin{questao}[3]{}
A fonte dependente vale $4i_x$ volts, com $i_x=I_1-I_2$ no resistor central.
Determine as correntes de malha.

\circuitoq{%
\begin{circuitikz}[american,scale=.92,transform shape,line width=.8pt]
\draw (0,0) to[\fonteV,l^={$\SI{24}{\volt}$}] (0,3.2)
      to[R,l={$\SI{4}{\ohm}$}] (3.2,3.2) coordinate(t);
\draw (t) to[R,l={$\SI{4}{\ohm}$},i>^={$i_x$}] (3.2,0) coordinate(b) -- (0,0);
\draw (t) to[R,l={$\SI{8}{\ohm}$}] (6.4,3.2)
      to[cvsource,l={$4i_x$}] (6.4,0) -- (b);
\draw[->,loopcol,line width=1.2pt] (1.75,1.00) arc[start angle=0,end angle=-310,radius=.46cm];
\node[loopcol,fill=white,inner sep=1pt] at (1.75,1.00){\footnotesize$I_1$};
\draw[->,loopcol,line width=1.2pt] (4.75,1.00) arc[start angle=0,end angle=-310,radius=.46cm];
\node[loopcol,fill=white,inner sep=1pt] at (4.75,1.00){\footnotesize$I_2$};
\end{circuitikz}}

\resposta{$I_1=\SI{4}{\ampere}$; $I_2=\SI{2}{\ampere}$}
\resolucao{A malha 1 fornece $8I_1-4I_2=24$. Na malha 2,
\[
-4I_1+12I_2=4(I_1-I_2),
\]
ou $-8I_1+16I_2=0$. Assim $I_1=2I_2$ e, pela primeira equação,
$I_1=\SI{4}{\ampere}$ e $I_2=\SI{2}{\ampere}$.}
\end{questao}

\begin{questao}[3]{}
A tensão de controle $v_x$ é medida no resistor de $\SI{4}{\ohm}$ exclusivo da
malha 1. A fonte dependente vale $0{,}5v_x$. Determine $I_1$ e $I_2$.

\circuitoq{%
\begin{circuitikz}[american,scale=.92,transform shape,line width=.8pt]
\draw (0,0) to[\fonteV,l^={$\SI{24}{\volt}$}] (0,3.2)
      to[R,l_={$\SI{4}{\ohm}$},v^>={$v_x$}] (3.2,3.2) coordinate(t);
\draw (t) to[R,l={$\SI{4}{\ohm}$}] (3.2,0) coordinate(b) -- (0,0);
\draw (t) to[R,l={$\SI{4}{\ohm}$}] (6.4,3.2)
      to[cvsource,l={$0{,}5v_x$}] (6.4,0) -- (b);
\draw[->,loopcol,line width=1.2pt] (1.75,1.00) arc[start angle=0,end angle=-310,radius=.46cm];
\node[loopcol,fill=white,inner sep=1pt] at (1.75,1.00){\footnotesize$I_1$};
\draw[->,loopcol,line width=1.2pt] (4.75,1.00) arc[start angle=0,end angle=-310,radius=.46cm];
\node[loopcol,fill=white,inner sep=1pt] at (4.75,1.00){\footnotesize$I_2$};
\end{circuitikz}}

\resposta{$I_1=\SI{4.8}{\ampere}$; $I_2=\SI{3.6}{\ampere}$}
\resolucao{Como $v_x=4I_1$, a fonte dependente vale $2I_1$ volts. Logo,
\[
8I_1-4I_2=24,
\qquad
-4I_1+8I_2=2I_1.
\]
Da segunda, $I_2=0{,}75I_1$. Substituindo na primeira, obtêm-se
$I_1=\SI{4.8}{\ampere}$ e $I_2=\SI{3.6}{\ampere}$.}
\end{questao}

\begin{questao}[3]{}
A fonte de corrente dependente está no ramo comum e fornece corrente de baixo
para cima igual a $0{,}1v_x$. A tensão $v_x$ está indicada no resistor da malha
1. Determine $I_1$ e $I_2$ por supermalha.

\circuitoq{%
\begin{circuitikz}[american,scale=.92,transform shape,line width=.8pt]
\draw (0,0) to[\fonteV,l^={$\SI{40}{\volt}$}] (0,3.2)
      to[R,l_={$\SI{5}{\ohm}$},v^>={$v_x$}] (3.2,3.2) coordinate(t);
\draw (t) to[cisource,l={$0{,}1v_x$},invert] (3.2,0) coordinate(b) -- (0,0);
\draw (t) to[R,l={$\SI{5}{\ohm}$}] (6.4,3.2) -- (6.4,0) -- (b);
\draw[->,loopcol,line width=1.2pt] (1.75,1.00) arc[start angle=0,end angle=-310,radius=.46cm];
\node[loopcol,fill=white,inner sep=1pt] at (1.75,1.00){\footnotesize$I_1$};
\draw[->,loopcol,line width=1.2pt] (4.75,1.00) arc[start angle=0,end angle=-310,radius=.46cm];
\node[loopcol,fill=white,inner sep=1pt] at (4.75,1.00){\footnotesize$I_2$};
\draw[dashed,laranja,line width=.9pt,rounded corners=7pt]
      (-.55,-.5) rectangle (6.95,4.10);
\node[laranja,font=\footnotesize\sffamily,fill=white,inner sep=1.5pt]
      at (3.2,-.80) {supermalha};
\end{circuitikz}}

\resposta{$I_1=\SI{3.2}{\ampere}$; $I_2=\SI{4.8}{\ampere}$}
\resolucao{Como $v_x=5I_1$, a restrição imposta pela fonte é
\[
I_2-I_1=0{,}1(5I_1)=0{,}5I_1,
\]
portanto $I_2=1{,}5I_1$. A LKT na supermalha externa é
$40=5I_1+5I_2$. Substituindo, $I_1=\SI{3.2}{\ampere}$ e
$I_2=\SI{4.8}{\ampere}$.}
\end{questao}

\begin{questao}[3]{}
A fonte de corrente dependente do ramo comum fornece, de baixo para cima,
$2i_x$, sendo $i_x=I_1$ no resistor esquerdo. Determine as correntes de malha
por supermalha.

\circuitoq{%
\begin{circuitikz}[american,scale=.92,transform shape,line width=.8pt]
\draw (0,0) to[\fonteV,l^={$\SI{40}{\volt}$}] (0,3.2)
      to[R,l_={$\SI{5}{\ohm}$},i>^={$i_x$}] (3.2,3.2) coordinate(t);
\draw (t) to[cisource,l={$2i_x$},invert] (3.2,0) coordinate(b) -- (0,0);
\draw (t) to[R,l={$\SI{5}{\ohm}$}] (6.4,3.2) -- (6.4,0) -- (b);
\draw[->,loopcol,line width=1.2pt] (1.75,1.00) arc[start angle=0,end angle=-310,radius=.46cm];
\node[loopcol,fill=white,inner sep=1pt] at (1.75,1.00){\footnotesize$I_1$};
\draw[->,loopcol,line width=1.2pt] (4.75,1.00) arc[start angle=0,end angle=-310,radius=.46cm];
\node[loopcol,fill=white,inner sep=1pt] at (4.75,1.00){\footnotesize$I_2$};
\draw[dashed,laranja,line width=.9pt,rounded corners=7pt]
      (-.55,-.5) rectangle (6.95,4.10);
\node[laranja,font=\footnotesize\sffamily,fill=white,inner sep=1.5pt]
      at (3.2,-.80) {supermalha};
\end{circuitikz}}

\resposta{$I_1=\SI{2}{\ampere}$; $I_2=\SI{6}{\ampere}$}
\resolucao{A restrição é $I_2-I_1=2I_1$, portanto $I_2=3I_1$. A LKT na
supermalha é $40=5I_1+5I_2$. Assim $40=20I_1$, resultando
$I_1=\SI{2}{\ampere}$ e $I_2=\SI{6}{\ampere}$.}
\end{questao}
%
    \fvMaxDesafioOrig
    % =====================================================================
% FONTES DEPENDENTES EM MALHAS --- NIVEL 4
% Questoes integradas ao bloco de Desafio do Capitulo 9.
% =====================================================================

\begin{questao}[4]{}
No circuito de três malhas, a fonte dependente da terceira malha vale
$2{,}5I_1$ volts e impulsiona $I_3$. Determine $I_1$, $I_2$ e $I_3$.

\circuitoq{%
\begin{circuitikz}[american,scale=.84,transform shape,line width=.8pt]
\draw (0,0) to[\fonteV,l^={$\SI{20}{\volt}$}] (0,3.2)
      to[R,l={$\SI{5}{\ohm}$}] (2.8,3.2) coordinate(t1);
\draw (t1) to[R,l_={$\SI{5}{\ohm}$}] (2.8,0) coordinate(b1) -- (0,0);
\draw (t1) to[R,l={$\SI{5}{\ohm}$}] (5.6,3.2) coordinate(t2);
\draw (t2) to[R,l_={$\SI{5}{\ohm}$}] (5.6,0) coordinate(b2) -- (b1);
\draw (t2) to[R,l={$\SI{5}{\ohm}$}] (8.4,3.2)
      to[cvsource,l={$2{,}5I_1$}] (8.4,0) -- (b2);
\foreach \x/\n in {1.25/1,4.15/2,7.05/3}{%
  \draw[->,loopcol,line width=1.05pt] (\x,1.00) arc (0:310:.40);
  \node[loopcol,fill=white,inner sep=1pt] at (\x,1.00){\scriptsize$I_{\n}$};}
\end{circuitikz}}

\resposta{$I_1=\dfrac{8}{3}\,\si{\ampere}$; $I_2=I_3=\dfrac{4}{3}\,\si{\ampere}$}
\resolucao{As três equações são
\[
10I_1-5I_2=20,
\]
\[
-5I_1+15I_2-5I_3=0,
\]
\[
-5I_2+10I_3=2{,}5I_1.
\]
A solução do sistema fornece $I_1=8/3$ A e $I_2=I_3=4/3$ A.}
\end{questao}

\begin{questao}[4]{}
A fonte de corrente dependente está entre as malhas 2 e 3 e impõe
$I_3-I_2=0{,}5I_1$. Determine as três correntes usando uma supermalha envolvendo
as malhas 2 e 3.

\circuitoq{%
\begin{circuitikz}[american,scale=.84,transform shape,line width=.8pt]
\draw (0,0) to[\fonteV,l^={$\SI{30}{\volt}$}] (0,3.2)
      to[R,l={$\SI{5}{\ohm}$}] (2.8,3.2) coordinate(t1);
\draw (t1) to[R,l_={$\SI{5}{\ohm}$}] (2.8,0) coordinate(b1) -- (0,0);
\draw (t1) to[R,l={$\SI{5}{\ohm}$}] (5.6,3.2) coordinate(t2);
\draw (t2) to[cisource,l={$0{,}5I_1$},invert] (5.6,0) coordinate(b2) -- (b1);
\draw (t2) to[R,l={$\SI{10}{\ohm}$}] (8.4,3.2) -- (8.4,0) -- (b2);
\foreach \x/\n in {1.25/1,4.15/2,7.05/3}{%
  \draw[->,loopcol,line width=1.05pt] (\x,1.00) arc (0:310:.40);
  \node[loopcol,fill=white,inner sep=1pt] at (\x,1.00){\scriptsize$I_{\n}$};}
\draw[dashed,laranja,line width=.9pt,rounded corners=7pt]
      (2.25,-.55) rectangle (8.95,4.05);
\node[laranja,font=\footnotesize\sffamily,fill=white,inner sep=1.5pt]
      at (5.6,-.85) {supermalha 2--3};
\end{circuitikz}}

\resposta{$I_1=\SI{3}{\ampere}$; $I_2=0$; $I_3=\SI{1.5}{\ampere}$}
\resolucao{A malha 1 fornece $10I_1-5I_2=30$. A supermalha 2--3 fornece
\[
-5I_1+10I_2+10I_3=0,
\]
e a fonte impõe $I_3-I_2=0{,}5I_1$. Resolvendo o sistema,
$I_1=\SI{3}{\ampere}$, $I_2=0$ e $I_3=\SI{1.5}{\ampere}$.}
\end{questao}

\begin{questao}[4]{}
A fonte dependente vale $kI_1$ volts. Determine $k$ para que
$I_2=\SI{2}{\ampere}$ e calcule $I_1$.

\circuitoq{%
\begin{circuitikz}[american,scale=.92,transform shape,line width=.8pt]
\draw (0,0) to[\fonteV,l^={$\SI{20}{\volt}$}] (0,3.2)
      to[R,l={$\SI{5}{\ohm}$}] (3.2,3.2) coordinate(t);
\draw (t) to[R,l={$\SI{5}{\ohm}$}] (3.2,0) coordinate(b) -- (0,0);
\draw (t) to[R,l={$\SI{5}{\ohm}$}] (6.4,3.2)
      to[cvsource,l={$kI_1$}] (6.4,0) -- (b);
\draw[->,loopcol,line width=1.2pt] (1.75,1.00) arc (0:310:.46);
\node[loopcol,fill=white,inner sep=1pt] at (1.75,1.00){\footnotesize$I_1$};
\draw[->,loopcol,line width=1.2pt] (4.75,1.00) arc (0:310:.46);
\node[loopcol,fill=white,inner sep=1pt] at (4.75,1.00){\footnotesize$I_2$};
\end{circuitikz}}

\resposta{$I_1=\SI{3}{\ampere}$; $k=\dfrac{5}{3}\,\ohm\approx\SI{1.67}{\ohm}$}
\resolucao{A malha 1 dá $10I_1-5I_2=20$. Impondo $I_2=2$ A,
$I_1=3$ A. Na malha 2,
\[
-5I_1+10I_2=kI_1.
\]
Logo $-15+20=3k$, e portanto $k=5/3\,\ohm$.}
\end{questao}

\begin{questao}[4]{}
A fonte de tensão dependente está no ramo comum e possui tensão, de cima para
baixo, $v_d=2I_1$. Determine $I_1$, $I_2$, a corrente real na fonte dependente e
sua potência, indicando se ela absorve ou fornece energia.

\circuitoq{%
\begin{circuitikz}[american,scale=.92,transform shape,line width=.8pt]
\draw (0,0) to[\fonteV,l^={$\SI{20}{\volt}$}] (0,3.2)
      to[R,l={$\SI{4}{\ohm}$}] (3.2,3.2) coordinate(t);
\draw (t) to[cvsource,l_={$v_d$}] (3.2,0) coordinate(b) -- (0,0);
\draw (t) to[R,l={$\SI{4}{\ohm}$}] (6.4,3.2) -- (6.4,0) -- (b);
\draw[->,loopcol,line width=1.2pt] (1.75,1.00) arc (0:310:.46);
\node[loopcol,fill=white,inner sep=1pt] at (1.75,1.00){\footnotesize$I_1$};
\draw[->,loopcol,line width=1.2pt] (4.75,1.00) arc (0:310:.46);
\node[loopcol,fill=white,inner sep=1pt] at (4.75,1.00){\footnotesize$I_2$};
\end{circuitikz}}

\resposta{$I_1=\dfrac{10}{3}\,\si{\ampere}$; $I_2=\dfrac{5}{3}\,\si{\ampere}$;
$i_d=\dfrac{5}{3}\,\si{\ampere}$ para baixo; $P_d=\dfrac{100}{9}\,\si{\watt}$ absorvidos}
\resolucao{Na malha 1, $4I_1+v_d=20$ e $v_d=2I_1$, portanto
$I_1=10/3$ A. Na malha 2, $4I_2-v_d=0$, então $I_2=5/3$ A. A corrente real no
ramo comum é $i_d=I_1-I_2=5/3$ A para baixo. Como $v_d=20/3$ V e a corrente
entra no terminal positivo superior da fonte dependente,
\[
P_d=v_di_d=\frac{20}{3}\frac{5}{3}=\frac{100}{9}\,\si{\watt},
\]
logo a fonte dependente \textbf{absorve} potência.}
\end{questao}
%
  \else
    \fvMaxDesafioOrig
  \fi
}

% =====================================================================
% BANCO COMPLEMENTAR --- ANALISE DE MALHAS / METODO DE MAXWELL
% Revisao 2026-08-15
% 7 N1 + 8 N2 + 5 N3 + 5 N4 = 25 questoes.
% N2/N3/N4: todas as questoes partem de CIRCUITOS diferentes e as setas
% de corrente de malha sao desenhadas explicitamente no sentido HORARIO.
% Fontes dependentes permanecem nos bancos especificos, evitando transformar
% toda a secao em repeticoes do mesmo tema.
% =====================================================================

\subsubsection*{Banco complementar --- Método de Maxwell (análise de malhas)}

\begin{atencao}[Convenção gráfica]
Em todos os circuitos deste banco, $I_1,I_2,\ldots$ são adotadas no
\textbf{sentido horário}. Uma corrente calculada negativa significa apenas que
a corrente real daquela malha circula no sentido oposto ao indicado. Os níveis
2, 3 e 4 usam topologias e objetivos diferentes: malhas triangulares, fontes em
ramos compartilhados, fontes de corrente periféricas, redes de três malhas,
ponte resistiva, síntese, máxima potência, identificação de parâmetros e rede
em roda de quatro malhas.
\end{atencao}

\providecommand{\fvMeshCW}[3]{%
  \draw[->,loopcol,line width=1.1pt] (#1)
    arc[start angle=0,end angle=-305,radius=#2cm];%
  \node[loopcol,fill=white,inner sep=1pt] at (#1){\footnotesize #3};}

% =====================================================================
% N1
% =====================================================================
\blocofixacao

\begin{questao}[1]{}
No desenho abaixo, a corrente de malha segue a convenção usada neste capítulo.
Identifique o sentido de $I_1$.
\circuitoq{%
\begin{circuitikz}[american,scale=.9,transform shape]
\draw (0,0) to[\fonteV,l^={$E$}] (0,2.8) to[R,l={$R$}] (4,2.8) -- (4,0)--(0,0);
\fvMeshCW{2.0,1.25}{.42}{$I_1$}
\end{circuitikz}}
\resposta{Horário}
\resolucao{A seta percorre o laço no sentido horário. Essa é apenas uma
convenção de referência; se o resultado algébrico de $I_1$ for negativo, a
corrente real circula no sentido anti-horário.}
\end{questao}

\begin{questao}[1]{}
Duas malhas horárias compartilham um resistor $R_c$. Qual é a corrente real no
resistor, tomando como positivo o sentido de $I_1$ no ramo comum?
\resposta{$i_c=I_1-I_2$}
\resolucao{As correntes horárias atravessam o ramo compartilhado em sentidos
opostos; por isso a corrente de ramo é a diferença algébrica.}
\end{questao}

\begin{questao}[1]{}
Um circuito planar conectado possui $b=11$ ramos e $n=7$ nós. Quantas correntes
de malha independentes são necessárias?
\resposta{$m=5$}
\resolucao{$m=b-n+1=11-7+1=5$.}
\end{questao}

\begin{questao}[1]{}
Uma fonte ideal de corrente está na periferia de uma única malha e sua seta
coincide com a corrente horária $I_2$. O que pode ser escrito imediatamente?
\resposta{$I_2=I_s$}
\resolucao{A fonte pertence somente àquela malha; portanto ela fixa diretamente
a corrente de malha e elimina uma incógnita do sistema.}
\end{questao}

\begin{questao}[1]{}
Uma fonte ideal de corrente está no ramo comum a duas malhas. Qual técnica deve
ser usada?
\resposta{Supermalha, mais a equação de restrição da fonte}
\resolucao{A tensão sobre uma fonte ideal de corrente é desconhecida; por isso
não se escreve LKT separadamente nas duas malhas. Contorna-se o ramo da fonte e
acrescenta-se a relação entre $I_1$ e $I_2$.}
\end{questao}

\begin{questao}[1]{}
Na matriz de uma rede resistiva com correntes de malha todas horárias, o que
representam $R_{kk}$ e $R_{jk}$, $j\ne k$?
\resposta{$R_{kk}$ é a soma das resistências da malha $k$; $R_{jk}$ é o negativo
da resistência compartilhada}
\resolucao{A diagonal reúne todos os resistores percorridos por $I_k$. Fora da
diagonal aparece $-R_c$ porque as correntes horárias percorrem o ramo comum em
sentidos opostos.}
\end{questao}

\begin{questao}[1]{}
Ao resolver uma malha adotada no sentido horário, obtém-se
$I_3=-\SI{0.8}{\ampere}$. Interprete fisicamente.
\resposta{A corrente real tem módulo $\SI{0.8}{\ampere}$ e circula no sentido anti-horário}
\resolucao{O sinal negativo não indica erro: ele informa que a escolha inicial
do sentido foi oposta à corrente física.}
\end{questao}

% =====================================================================
% N2 --- circuitos distintos
% =====================================================================
\blocoaplicacao

\begin{questao}[2]{}
Duas malhas triangulares compartilham o resistor diagonal de
$\SI{6}{\ohm}$. Determine $I_1$, $I_2$ e a corrente no ramo compartilhado.
\circuitoq{%
\begin{circuitikz}[american,scale=.88,transform shape,line width=.8pt]
\coordinate (A) at (0,0); \coordinate (B) at (3.4,4.2); \coordinate (C) at (6.8,0);
\draw (A) to[\fonteV,l^={$\SI{18}{\volt}$}] (0,2.0)
      to[R,l={$\SI{4}{\ohm}$}] (B);
\draw (B) to[R,l={$\SI{3}{\ohm}$}] (5.2,2.1)
      to[R,l={$\SI{3}{\ohm}$}] (C);
\draw (C)--(A);
\draw (B) to[R,l={$\SI{6}{\ohm}$}] (3.4,0)--(A);
\fvMeshCW{1.75,1.45}{.38}{$I_1$}
\fvMeshCW{5.05,1.45}{.38}{$I_2$}
\end{circuitikz}}
\resposta{$I_1=\dfrac{18}{7}\,\si{\ampere}$; $I_2=\dfrac{9}{7}\,\si{\ampere}$;
$i_c=\dfrac{9}{7}\,\si{\ampere}$}
\resolucao{As equações são
\[
10I_1-6I_2=18,\qquad -6I_1+12I_2=0.
\]
Da segunda, $I_1=2I_2$. Logo $14I_2=18$, resultando
$I_2=9/7$ A e $I_1=18/7$ A. No ramo comum,
$i_c=I_1-I_2=9/7$ A.}
\end{questao}

\begin{questao}[2]{}
A fonte de corrente de $\SI{1.5}{\ampere}$ está na periferia da malha 2 e sua
seta coincide com $I_2$. Determine $I_1$, $I_2$ e a corrente no resistor
central de $\SI{5}{\ohm}$.
\circuitoq{%
\begin{circuitikz}[american,scale=.92,transform shape,line width=.8pt]
\draw (0,0) to[\fonteV,l^={$\SI{30}{\volt}$}] (0,3.2)
      to[R,l={$\SI{5}{\ohm}$}] (3.2,3.2) coordinate(t);
\draw (t) to[R,l_={$\SI{5}{\ohm}$}] (3.2,0) coordinate(b)--(0,0);
\draw (t) to[R,l={$\SI{10}{\ohm}$}] (6.4,3.2)
      to[isource,l={$\SI{1.5}{\ampere}$}] (6.4,0)--(b);
\fvMeshCW{1.55,1.05}{.40}{$I_1$}
\fvMeshCW{4.85,1.05}{.40}{$I_2$}
\end{circuitikz}}
\resposta{$I_1=\SI{3.75}{\ampere}$; $I_2=\SI{1.5}{\ampere}$;
$i_c=\SI{2.25}{\ampere}$}
\resolucao{A fonte periférica fixa $I_2=1{,}5$ A. Para a malha 1,
\[
5I_1+5(I_1-I_2)=30,
\]
portanto $10I_1-7{,}5=30$ e $I_1=3{,}75$ A. Assim
$i_c=I_1-I_2=2{,}25$ A.}
\end{questao}

\begin{questao}[2]{}
O ramo comum contém um resistor de $\SI{4}{\ohm}$ em série com uma fonte de
$\SI{6}{\volt}$, com terminal positivo no topo. Determine $I_1$ e $I_2$.
\circuitoq{%
\begin{circuitikz}[american,scale=.9,transform shape,line width=.8pt]
\draw (0,0) to[\fonteV,l^={$\SI{24}{\volt}$}] (0,3.6)
      to[R,l={$\SI{4}{\ohm}$}] (3.4,3.6) coordinate(t);
\draw (t) to[R,l_={$\SI{4}{\ohm}$}] (3.4,2.0)
      to[\fonteV,l_={$\SI{6}{\volt}$}] (3.4,0) coordinate(b)--(0,0);
\draw (t) to[R,l={$\SI{8}{\ohm}$}] (6.8,3.6)--(6.8,0)--(b);
\fvMeshCW{1.65,1.15}{.40}{$I_1$}
\fvMeshCW{5.15,1.15}{.40}{$I_2$}
\end{circuitikz}}
\resposta{$I_1=\SI{3}{\ampere}$; $I_2=\SI{1.5}{\ampere}$}
\resolucao{Percorrendo as duas malhas no sentido horário:
\[
8I_1-4I_2=18,
\qquad
-4I_1+12I_2=6.
\]
A solução é $I_1=3$ A e $I_2=1{,}5$ A. A fonte compartilhada entra com sinais
opostos nas duas equações porque é atravessada em sentidos opostos.}
\end{questao}

\begin{questao}[2]{}
A rede possui três malhas triangulares em torno do nó central $O$. Monte a
matriz por inspeção e determine as três correntes.
\circuitoq{%
\begin{circuitikz}[american,scale=.78,transform shape,line width=.8pt]
\coordinate (A) at (0,0); \coordinate (B) at (4,6); \coordinate (C) at (8,0); \coordinate (O) at (4,2.1);
% borda AB: R=4 e E1=16
\draw (A) to[\fonteV,l^={$\SI{16}{\volt}$}] (1.15,1.75)
      to[R,l={$\SI{4}{\ohm}$}] (B);
% borda BC: R=5 e E2=11
\draw (B) to[R,l={$\SI{5}{\ohm}$}] (6.85,1.75)
      to[\fonteV,l={$\SI{11}{\volt}$}] (C);
% borda CA: R=6 e E3=1
\draw (C) to[R,l_={$\SI{6}{\ohm}$}] (4.7,0)
      to[\fonteV,l_={$\SI{1}{\volt}$}] (A);
% raios compartilhados
\draw (B) to[R,l={$\SI{2}{\ohm}$}] (O);
\draw (C) to[R,l={$\SI{3}{\ohm}$}] (O);
\draw (A) to[R,l_={$\SI{1}{\ohm}$}] (O);
\fvMeshCW{2.8,3.5}{.34}{$I_1$}
\fvMeshCW{5.4,3.5}{.34}{$I_2$}
\fvMeshCW{4.0,.95}{.34}{$I_3$}
\node[fill=white,inner sep=1pt] at(O){$O$};
\end{circuitikz}}
\resposta{$I_1=\SI{3}{\ampere}$; $I_2=\SI{2}{\ampere}$; $I_3=\SI{1}{\ampere}$}
\resolucao{Por inspeção,
\[
\begin{bmatrix}
7&-2&-1\\
-2&10&-3\\
-1&-3&10
\end{bmatrix}
\begin{bmatrix}I_1\\I_2\\I_3\end{bmatrix}
=
\begin{bmatrix}16\\11\\1\end{bmatrix}.
\]
Substituindo $(3,2,1)$ verifica-se cada linha; logo essa é a solução. O desenho
mostra que a matriz depende de quais raios são compartilhados, e não de uma
sequência linear de malhas.}
\end{questao}

\begin{questao}[2]{}
As duas fontes do circuito atuam em oposição. Determine as correntes horárias e
interprete o sinal de $I_2$.
\circuitoq{%
\begin{circuitikz}[american,scale=.92,transform shape,line width=.8pt]
\draw (0,0) to[\fonteV,l^={$\SI{20}{\volt}$}] (0,3.2)
      to[R,l={$\SI{6}{\ohm}$}] (3.2,3.2) coordinate(t);
\draw (t) to[R,l_={$\SI{4}{\ohm}$}] (3.2,0) coordinate(b)--(0,0);
\draw (t) to[R,l={$\SI{6}{\ohm}$}] (6.4,3.2)
      to[\fonteV,l_={$\SI{12}{\volt}$}] (6.4,0)--(b);
\fvMeshCW{1.55,1.05}{.40}{$I_1$}
\fvMeshCW{4.85,1.05}{.40}{$I_2$}
\end{circuitikz}}
\resposta{$I_1=\dfrac{38}{21}\,\si{\ampere}$;
$I_2=-\dfrac{10}{21}\,\si{\ampere}$}
\resolucao{O sistema é
\[
\begin{bmatrix}10&-4\\-4&10\end{bmatrix}
\begin{bmatrix}I_1\\I_2\end{bmatrix}
=
\begin{bmatrix}20\\-12\end{bmatrix}.
\]
Daí $I_1=38/21$ A e $I_2=-10/21$ A. Portanto a corrente física da segunda
malha tem módulo $10/21$ A e circula no sentido anti-horário.}
\end{questao}

\begin{questao}[2]{}
Escolha o valor e a polaridade da fonte $E$ para que a corrente no resistor
compartilhado de $\SI{4}{\ohm}$ seja nula.
\circuitoq{%
\begin{circuitikz}[american,scale=.92,transform shape,line width=.8pt]
\draw (0,0) to[\fonteV,l^={$\SI{24}{\volt}$}] (0,3.2)
      to[R,l={$\SI{4}{\ohm}$}] (3.2,3.2) coordinate(t);
\draw (t) to[R,l_={$\SI{4}{\ohm}$}] (3.2,0) coordinate(b)--(0,0);
\draw (t) to[R,l={$\SI{8}{\ohm}$}] (6.4,3.2)
      to[\fonteV,l_={$E$}] (6.4,0)--(b);
\fvMeshCW{1.55,1.05}{.40}{$I_1$}
\fvMeshCW{4.85,1.05}{.40}{$I_2$}
\end{circuitikz}}
\resposta{$E=\SI{48}{\volt}$, impulsionando $I_2$ no sentido horário}
\resolucao{Corrente nula no ramo comum exige $I_1=I_2=I$. A primeira malha dá
$(8-4)I=24$, ou $4I=24$, logo $I=6$ A. Na segunda,
$(12-4)I=E$, portanto $E=8\cdot6=48$ V, com polaridade que impulsione $I_2$.}
\end{questao}

\begin{questao}[2]{}
Determine as correntes de malha e faça o balanço de potência.
\circuitoq{%
\begin{circuitikz}[american,scale=.92,transform shape,line width=.8pt]
\draw (0,0) to[\fonteV,l^={$\SI{36}{\volt}$}] (0,3.2)
      to[R,l={$\SI{6}{\ohm}$}] (3.2,3.2) coordinate(t);
\draw (t) to[R,l_={$\SI{3}{\ohm}$}] (3.2,0) coordinate(b)--(0,0);
\draw (t) to[R,l={$\SI{6}{\ohm}$}] (6.4,3.2)
      to[\fonteV,l_={$\SI{12}{\volt}$},invert] (6.4,0)--(b);
\fvMeshCW{1.55,1.05}{.40}{$I_1$}
\fvMeshCW{4.85,1.05}{.40}{$I_2$}
\end{circuitikz}}
\resposta{$I_1=\SI{5}{\ampere}$; $I_2=\SI{3}{\ampere}$;
$P_{fontes}=P_R=\SI{216}{\watt}$}
\resolucao{O sistema
\[
\begin{bmatrix}9&-3\\-3&9\end{bmatrix}
\begin{bmatrix}I_1\\I_2\end{bmatrix}
=
\begin{bmatrix}36\\12\end{bmatrix}
\]
fornece $(I_1,I_2)=(5,3)$ A. A corrente no resistor comum é $2$ A. Assim
\[
P_R=6(5)^2+6(3)^2+3(2)^2=150+54+12=216\,\text{W}.
\]
As fontes fornecem $36(5)+12(3)=216$ W.}
\end{questao}

\begin{questao}[2]{}
A fonte de corrente da malha 2 fixa $I_2=\SI{2}{\ampere}$. Determine $I_1$ e a
tensão $v_s$ da fonte de corrente, definida positiva do terminal superior para
o inferior.
\circuitoq{%
\begin{circuitikz}[american,scale=.92,transform shape,line width=.8pt]
\draw (0,0) to[\fonteV,l^={$\SI{20}{\volt}$}] (0,3.2)
      to[R,l={$\SI{5}{\ohm}$}] (3.2,3.2) coordinate(t);
\draw (t) to[R,l_={$\SI{5}{\ohm}$}] (3.2,0) coordinate(b)--(0,0);
\draw (t) to[R,l={$\SI{10}{\ohm}$}] (6.4,3.2)
      to[isource,l={$\SI{2}{\ampere}$}] (6.4,0)--(b);
\node[right] at(6.75,1.6){$v_s$};
\fvMeshCW{1.55,1.05}{.40}{$I_1$}
\fvMeshCW{4.85,1.05}{.40}{$I_2$}
\end{circuitikz}}
\resposta{$I_1=\SI{3}{\ampere}$; $v_s=\SI{-15}{\volt}$}
\resolucao{Na malha 1,
$5I_1+5(I_1-2)=20$, então $I_1=3$ A. Na malha 2, percorrida no sentido horário,
\[
10(2)+v_s-5(3-2)=0,
\]
portanto $v_s=-15$ V. Isso significa que o terminal inferior da fonte está
$15$ V acima do terminal superior.}
\end{questao}

% =====================================================================
% N3 --- aprofundamento com novas topologias/objetivos
% =====================================================================
\blocoaprofundamento

\begin{questao}[3]{}
A fonte de corrente de $\SI{2}{\ampere}$ está entre as malhas 1 e 2. A fonte de
$\SI{8}{\volt}$ da malha 3 se opõe ao sentido horário indicado. Resolva usando
uma supermalha 1--2.
\circuitoq{%
\begin{circuitikz}[american,scale=.82,transform shape,line width=.8pt]
\draw (0,0) to[\fonteV,l^={$\SI{24}{\volt}$}] (0,3.4)
      to[R,l={$\SI{4}{\ohm}$}] (2.8,3.4) coordinate(t1);
\draw (t1) to[isource,l={$\SI{2}{\ampere}$},invert] (2.8,0) coordinate(b1)--(0,0);
\draw (t1) to[R,l={$\SI{2}{\ohm}$}] (5.6,3.4) coordinate(t2);
\draw (t2) to[R,l_={$\SI{4}{\ohm}$}] (5.6,0) coordinate(b2)--(b1);
\draw (t2) to[R,l={$\SI{4}{\ohm}$}] (8.4,3.4)
      to[\fonteV,l_={$\SI{8}{\volt}$}] (8.4,0)--(b2);
\fvMeshCW{1.25,1.05}{.34}{$I_1$}
\fvMeshCW{4.15,1.05}{.34}{$I_2$}
\fvMeshCW{7.05,1.05}{.34}{$I_3$}
\draw[dashed,laranja,line width=.9pt,rounded corners=7pt] (-.45,-.45) rectangle (6.05,3.85);
\node[laranja,fill=white,inner sep=1pt] at(2.8,-.72){\footnotesize supermalha 1--2};
\end{circuitikz}}
\resposta{$I_1=\SI{1.5}{\ampere}$; $I_2=\SI{3.5}{\ampere}$;
$I_3=\SI{0.75}{\ampere}$}
\resolucao{As equações são
\[
4I_1+6I_2-4I_3=24,
\qquad I_2-I_1=2,
\qquad -4I_2+8I_3=-8.
\]
Resolvendo, obtém-se $I_1=3/2$ A, $I_2=7/2$ A e $I_3=3/4$ A.}
\end{questao}

\begin{questao}[3]{}
O ramo comum é uma fonte ideal de $\SI{6}{\volt}$, positiva no topo. A fonte de
$\SI{18}{\volt}$ impulsiona $I_1$, enquanto a de $\SI{12}{\volt}$ se opõe a
$I_2$. Determine as correntes e a potência da fonte compartilhada.
\circuitoq{%
\begin{circuitikz}[american,scale=.92,transform shape,line width=.8pt]
\draw (0,0) to[\fonteV,l^={$\SI{18}{\volt}$}] (0,3.2)
      to[R,l={$\SI{6}{\ohm}$}] (3.2,3.2) coordinate(t);
\draw (t) to[\fonteV,l_={$\SI{6}{\volt}$}] (3.2,0) coordinate(b)--(0,0);
\draw (t) to[R,l={$\SI{4}{\ohm}$}] (6.4,3.2)
      to[\fonteV,l_={$\SI{12}{\volt}$}] (6.4,0)--(b);
\fvMeshCW{1.55,1.05}{.40}{$I_1$}
\fvMeshCW{4.85,1.05}{.40}{$I_2$}
\end{circuitikz}}
\resposta{$I_1=\SI{2}{\ampere}$; $I_2=\SI{-1.5}{\ampere}$;
$P_{6V}=\SI{21}{\watt}$ absorvidos}
\resolucao{Na malha 1, $6I_1+6=18$, logo $I_1=2$ A. Na malha 2,
$4I_2+12-6=0$, então $I_2=-1{,}5$ A. A corrente real na fonte comum, para
baixo, é $I_1-I_2=3{,}5$ A. Como ela entra no terminal positivo superior,
\[
P=6(3{,}5)=21\,\text{W},
\]
portanto a fonte de $6$ V absorve potência.}
\end{questao}

\begin{questao}[3]{}
As fontes de corrente periféricas fixam $I_1=\SI{2}{\ampere}$ e
$I_3=\SI{1}{\ampere}$. Determine $I_2$ e identifique qualquer ramo
compartilhado com corrente nula.
\circuitoq{%
\begin{circuitikz}[american,scale=.82,transform shape,line width=.8pt]
\draw (0,0) to[isource,l^={$\SI{2}{\ampere}$}] (0,3.4)
      --(2.8,3.4) coordinate(t1);
\draw (t1) to[R,l_={$\SI{2}{\ohm}$}] (2.8,0) coordinate(b1)--(0,0);
\draw (t1) to[R,l={$\SI{4}{\ohm}$}] (4.5,3.4)
      to[\fonteV,l={$\SI{11}{\volt}$}] (5.8,3.4) coordinate(t2);
\draw (t2) to[R,l_={$\SI{3}{\ohm}$}] (5.8,0) coordinate(b2)--(b1);
\draw (t2)--(8.6,3.4) to[isource,l={$\SI{1}{\ampere}$}] (8.6,0)--(b2);
\fvMeshCW{1.25,1.05}{.34}{$I_1$}
\fvMeshCW{4.30,1.05}{.34}{$I_2$}
\fvMeshCW{7.25,1.05}{.34}{$I_3$}
\end{circuitikz}}
\resposta{$I_2=\SI{2}{\ampere}$; o resistor de $\SI{2}{\ohm}$ conduz corrente nula}
\resolucao{A equação da malha central é
\[
(2+4+3)I_2-2I_1-3I_3=11.
\]
Com $I_1=2$ A e $I_3=1$ A:
$9I_2-4-3=11$, portanto $I_2=2$ A. Assim, no resistor compartilhado entre as
malhas 1 e 2, $I_1-I_2=0$.}
\end{questao}

\begin{questao}[3]{}
Analise a ponte resistiva pelo método das malhas e determine a corrente no
resistor de ponte de $\SI{5}{\ohm}$.
\circuitoq{%
\begin{circuitikz}[american,scale=.78,transform shape,line width=.8pt]
\coordinate (A) at (5,5); \coordinate (D) at (5,0); \coordinate (B) at (2.6,2.5); \coordinate (C) at (7.4,2.5);
\draw (A)--(0.5,5) to[\fonteV,l^={$\SI{18}{\volt}$}] (0.5,0)--(D);
\draw (A) to[R,l_={$\SI{2}{\ohm}$}] (B) to[R,l_={$\SI{4}{\ohm}$}] (D);
\draw (A) to[R,l={$\SI{3}{\ohm}$}] (C) to[R,l={$\SI{6}{\ohm}$}] (D);
\draw (B) to[R,l={$\SI{5}{\ohm}$}] (C);
\fvMeshCW{2.1,2.55}{.34}{$I_1$}
\fvMeshCW{5.0,3.55}{.32}{$I_2$}
\fvMeshCW{5.0,1.45}{.32}{$I_3$}
\end{circuitikz}}
\resposta{$I_1=\SI{5}{\ampere}$; $I_2=I_3=\SI{2}{\ampere}$;
$i_{ponte}=0$}
\resolucao{A matriz é
\[
\begin{bmatrix}
6&-2&-4\\
-2&10&-5\\
-4&-5&15
\end{bmatrix}
\begin{bmatrix}I_1\\I_2\\I_3\end{bmatrix}
=
\begin{bmatrix}18\\0\\0\end{bmatrix}.
\]
A solução é $(5,2,2)$ A. O resistor de ponte é comum às malhas 2 e 3, logo
$i_{ponte}=I_2-I_3=0$. O resultado expressa o equilíbrio da ponte.}
\end{questao}

\begin{questao}[3]{}
Use uma fonte de teste de $\SI{12}{\volt}$ e análise de malhas para determinar
a resistência equivalente vista pela fonte.
\circuitoq{%
\begin{circuitikz}[american,scale=.92,transform shape,line width=.8pt]
\draw (0,0) to[\fonteV,l^={$\SI{12}{\volt}$}] (0,3.2)
      to[R,l={$\SI{3}{\ohm}$}] (3.2,3.2) coordinate(t);
\draw (t) to[R,l_={$\SI{6}{\ohm}$}] (3.2,0) coordinate(b)--(0,0);
\draw (t) to[R,l={$\SI{6}{\ohm}$}] (6.4,3.2)--(6.4,0)--(b);
\fvMeshCW{1.55,1.05}{.40}{$I_1$}
\fvMeshCW{4.85,1.05}{.40}{$I_2$}
\end{circuitikz}}
\resposta{$I_1=\SI{2}{\ampere}$; $I_2=\SI{1}{\ampere}$;
$R_{eq}=\SI{6}{\ohm}$}
\resolucao{As equações são
\[
9I_1-6I_2=12,\qquad -6I_1+12I_2=0.
\]
Da segunda, $I_1=2I_2$; substituindo na primeira, $I_2=1$ A e $I_1=2$ A. A
fonte de teste fornece $2$ A, portanto $R_{eq}=12/2=6\,\ohm$.}
\end{questao}

% =====================================================================
% N4 --- sintese, otimizacao e problemas inversos
% =====================================================================
\blocodesafio

\begin{questao}[4]{}
A rede em ``roda'' possui quatro malhas triangulares em torno do nó $O$.
Monte a matriz e determine as quatro correntes.
\circuitoq{%
\begin{circuitikz}[american,scale=.72,transform shape,line width=.8pt]
\coordinate (A) at (-4,4); \coordinate (B) at (4,4); \coordinate (C) at (4,-4); \coordinate (D) at (-4,-4); \coordinate (O) at (0,0);
% lados externos: R=4 em serie com fontes
\draw (A) to[\fonteV,l^={$\SI{24}{\volt}$}] (-1.8,4) to[R,l={$\SI{4}{\ohm}$}] (B);
\draw (B) to[\fonteV,l={$\SI{12}{\volt}$}] (4,1.8) to[R,l={$\SI{4}{\ohm}$}] (C);
\draw (C) to[\fonteV,l_={$\SI{8}{\volt}$}] (1.8,-4) to[R,l_={$\SI{4}{\ohm}$}] (D);
\draw (D) to[\fonteV,l_={$\SI{4}{\volt}$},invert] (-4,1.8) to[R,l_={$\SI{4}{\ohm}$}] (A);
% raios
\draw (A) to[R,l={$\SI{2}{\ohm}$}] (O);
\draw (B) to[R,l_={$\SI{2}{\ohm}$}] (O);
\draw (C) to[R,l={$\SI{2}{\ohm}$}] (O);
\draw (D) to[R,l_={$\SI{2}{\ohm}$}] (O);
\fvMeshCW{0,2.45}{.34}{$I_1$}
\fvMeshCW{2.45,0}{.34}{$I_2$}
\fvMeshCW{0,-2.45}{.34}{$I_3$}
\fvMeshCW{-2.45,0}{.34}{$I_4$}
\node[fill=white,inner sep=1pt] at(O){$O$};
\end{circuitikz}}
\resposta{$I_1=\SI{4}{\ampere}$; $I_2=\SI{3}{\ampere}$;
$I_3=\SI{2}{\ampere}$; $I_4=\SI{1}{\ampere}$}
\resolucao{Cada malha contém $4\,\ohm$ exclusivo e dois raios de
$2\,\ohm$. A matriz é cíclica:
\[
\begin{bmatrix}
8&-2&0&-2\\
-2&8&-2&0\\
0&-2&8&-2\\
-2&0&-2&8
\end{bmatrix}
\begin{bmatrix}I_1\\I_2\\I_3\\I_4\end{bmatrix}
=
\begin{bmatrix}24\\12\\8\\-4\end{bmatrix}.
\]
A solução é $(4,3,2,1)$ A. O termo $-4$ representa a fonte da quarta malha
orientada contra a corrente horária de referência.}
\end{questao}

\begin{questao}[4]{}
Determine o resistor compartilhado $R$ para que $I_2=I_1/2$. Calcule também
as correntes resultantes.
\circuitoq{%
\begin{circuitikz}[american,scale=.92,transform shape,line width=.8pt]
\draw (0,0) to[\fonteV,l^={$\SI{20}{\volt}$}] (0,3.2)
      to[R,l={$\SI{4}{\ohm}$}] (3.2,3.2) coordinate(t);
\draw (t) to[R,l_={$R$}] (3.2,0) coordinate(b)--(0,0);
\draw (t) to[R,l={$\SI{6}{\ohm}$}] (6.4,3.2)--(6.4,0)--(b);
\fvMeshCW{1.55,1.05}{.40}{$I_1$}
\fvMeshCW{4.85,1.05}{.40}{$I_2$}
\end{circuitikz}}
\resposta{$R=\SI{6}{\ohm}$; $I_1=\dfrac{20}{7}\,\si{\ampere}$;
$I_2=\dfrac{10}{7}\,\si{\ampere}$}
\resolucao{Na malha 2,
\[
-RI_1+(6+R)I_2=0.
\]
Impondo $I_2=I_1/2$:
$-R+(6+R)/2=0$, logo $R=6\,\ohm$. Na malha 1,
$10I_1-6I_2=20$; com $I_2=I_1/2$, resulta $7I_1=20$.}
\end{questao}

\begin{questao}[4]{}
O resistor compartilhado é a carga $R_L$. Determine $R_L$ para que sua potência
seja máxima e calcule $P_{L,\max}$.
\circuitoq{%
\begin{circuitikz}[american,scale=.92,transform shape,line width=.8pt]
\draw (0,0) to[\fonteV,l^={$\SI{24}{\volt}$}] (0,3.2)
      to[R,l={$\SI{4}{\ohm}$}] (3.2,3.2) coordinate(t);
\draw (t) to[R,l_={$R_L$}] (3.2,0) coordinate(b)--(0,0);
\draw (t) to[R,l={$\SI{8}{\ohm}$}] (6.4,3.2)
      to[\fonteV,l_={$\SI{12}{\volt}$},invert] (6.4,0)--(b);
\fvMeshCW{1.55,1.05}{.40}{$I_1$}
\fvMeshCW{4.85,1.05}{.40}{$I_2$}
\end{circuitikz}}
\resposta{$R_L=\dfrac{8}{3}\,\ohm\approx\SI{2.67}{\ohm}$;
$P_{L,\max}=\SI{13.5}{\watt}$}
\resolucao{As equações são
\[
(4+R_L)I_1-R_LI_2=24,
\]
\[
-R_LI_1+(8+R_L)I_2=12.
\]
A corrente na carga é
\[
i_L=I_1-I_2=\frac{36}{3R_L+8}.
\]
Logo
\[
P_L=R_Li_L^2=\frac{1296R_L}{(3R_L+8)^2}.
\]
Derivando e igualando a zero, obtém-se $R_L=8/3\,\ohm$. Nesse ponto,
$i_L=36/16=2{,}25$ A e $P_{L,\max}=13{,}5$ W.}
\end{questao}

\begin{questao}[4]{}
A fonte $E$ deve ser escolhida para impor $I_2=0$. Determine $E$, sua polaridade
relativa ao sentido horário e as correntes $I_1$ e $I_3$.
\circuitoq{%
\begin{circuitikz}[american,scale=.82,transform shape,line width=.8pt]
\draw (0,0) to[\fonteV,l^={$\SI{24}{\volt}$}] (0,3.4)
      to[R,l={$\SI{6}{\ohm}$}] (2.8,3.4) coordinate(t1);
\draw (t1) to[R,l_={$\SI{2}{\ohm}$}] (2.8,0) coordinate(b1)--(0,0);
\draw (t1) to[R,l={$\SI{4}{\ohm}$}] (5.8,3.4) coordinate(t2);
\draw (t2) to[R,l_={$\SI{1}{\ohm}$}] (5.8,0) coordinate(b2)--(b1);
\draw (t2) to[R,l={$\SI{4}{\ohm}$}] (8.8,3.4)
      to[\fonteV,l_={$E$}] (8.8,0)--(b2);
\fvMeshCW{1.25,1.05}{.34}{$I_1$}
\fvMeshCW{4.30,1.05}{.34}{$I_2$}
\fvMeshCW{7.35,1.05}{.34}{$I_3$}
\end{circuitikz}}
\resposta{$I_1=\SI{3}{\ampere}$; $I_2=0$; $I_3=\SI{-6}{\ampere}$;
$E=\SI{30}{\volt}$ orientada contra $I_3$ horário}
\resolucao{O sistema é
\[
8I_1-2I_2=24,
\qquad
-2I_1+7I_2-I_3=0,
\qquad
-I_2+5I_3=E_s,
\]
onde $E_s$ é positivo se impulsiona $I_3$ horário. Impondo $I_2=0$:
$I_1=3$ A e, pela segunda equação, $I_3=-6$ A. Assim
$E_s=5(-6)=-30$ V. Portanto a fonte deve ter magnitude $30$ V e polaridade
oposta à referência horária.}
\end{questao}

\begin{questao}[4]{}
Em ensaio de uma rede de duas malhas, mediram-se $I_1=\SI{3}{\ampere}$ e
$I_2=\SI{2}{\ampere}$. Determine a resistência desconhecida $R$ do ramo comum
e a potência dissipada nela.
\circuitoq{%
\begin{circuitikz}[american,scale=.92,transform shape,line width=.8pt]
\draw (0,0) to[\fonteV,l^={$\SI{20}{\volt}$}] (0,3.2)
      to[R,l={$\SI{4}{\ohm}$}] (3.2,3.2) coordinate(t);
\draw (t) to[R,l_={$R$}] (3.2,0) coordinate(b)--(0,0);
\draw (t) to[R,l={$\SI{6}{\ohm}$}] (6.4,3.2)
      to[\fonteV,l_={$\SI{4}{\volt}$},invert] (6.4,0)--(b);
\fvMeshCW{1.55,1.05}{.40}{$I_1$}
\fvMeshCW{4.85,1.05}{.40}{$I_2$}
\end{circuitikz}}
\resposta{$R=\SI{8}{\ohm}$; $P_R=\SI{8}{\watt}$}
\resolucao{Pela malha 1,
\[
(4+R)3-R(2)=20 \Rightarrow 12+R=20 \Rightarrow R=8\,\ohm.
\]
A malha 2 confirma o mesmo valor:
\[
-3R+(6+R)2=4 \Rightarrow 12-R=4 \Rightarrow R=8\,\ohm.
\]
A corrente no ramo comum é $I_1-I_2=1$ A; portanto
$P_R=8(1)^2=8$ W.}
\end{questao}

\endgroup

\gabarito[Capítulo 9 --- Método de Maxwell]
